% !TeX root = ../初等数学笔记.tex
% 由主文档通过 \input 载入；请编译主文档。
\setcounter{chapter}{21} % 保留原书章号。
\chapter{有限概率模型}
\label{chap:em-probability}
概率为随机现象指定可计算的模型。随机并不意味着每种描述都等可能，概率计算必须先说明基本结果是什么，以及为什么采用所给权重。本章只讨论有限样本空间，避免引入尚未建立的测度理论。

\section{随机试验、样本空间与事件}
\emph{样本空间}（sample space）$\Omega$ 是一次试验所有可能基本结果组成的有限集合，\emph{事件}（event）是其子集。两次掷硬币可用有序对 $(\text{正},\text{反})$ 等表示；“至少一次正面”包含三个有序结果，而“正面次数为零、一、二”这三类大小不同，不能因为有三类就各赋概率 $1/3$。

并、交、补分别表达至少一个发生、同时发生和不发生。$A\cap B=\varnothing$ 称为互斥，$B=A^c$ 称为对立；对立还要求并为整个样本空间，因此比互斥更强。基本结果的区分应与试验记录和概率赋值相匹配。

\section{概率、古典概型与加法公式}
给每个 $\omega\in\Omega$ 指定 $p(\omega)\ge0$ 且 $\sum_\omega p(\omega)=1$，定义
\[
 \mathbb P(A)=\sum_{\omega\in A}p(\omega).
\]
于是 $\mathbb P(\varnothing)=0,\mathbb P(\Omega)=1$，互斥事件概率可相加。有限等可能模型才满足 $\mathbb P(A)=|A|/|\Omega|$。对一般事件，分割为不重叠部分得到
\[
 \mathbb P(A\cup B)=\mathbb P(A)+\mathbb P(B)-\mathbb P(A\cap B),
 \qquad \mathbb P(A^c)=1-\mathbb P(A).
\]
例如从 $5$ 个不同球中等可能无放回抽取 $2$ 个，若其中 $2$ 个为红球，则两红概率为 $\binom22/\binom52=1/10$。采用无序结果是合法的，因为每个无序对对应相同数量的有序抽取。

\section{条件概率、独立与互斥}
在 $\mathbb P(B)>0$ 时，定义\emph{条件概率}（conditional probability）
\[
 \mathbb P(A\mid B)=\frac{\mathbb P(A\cap B)}{\mathbb P(B)}.
\]
它把注意范围缩到 $B$，再重新归一化。由定义得乘法公式 $\mathbb P(A\cap B)=\mathbb P(B)\mathbb P(A\mid B)$，并可逐步扩展到多次条件选择。

事件 $A,B$ \emph{独立}（independent）是指 $\mathbb P(A\cap B)=\mathbb P(A)\mathbb P(B)$；当 $\mathbb P(B)>0$ 时等价于条件信息不改变 $A$ 的概率。两个正概率互斥事件不独立，因为交集概率为零而概率乘积为正。多个事件联合独立要求每个子族的交概率都等于相应乘积，仅两两独立不充分。例如两个独立公平比特为 $X,Y$，事件 $X=0,Y=0,X=Y$ 两两独立，但三者同时发生概率为 $1/4$，不等于 $(1/2)^3$。

\section{全概率与贝叶斯公式}
设 $B_1,\ldots,B_m$ 两两互斥且并为 $\Omega$，并有 $\mathbb P(B_i)>0$。按来源分割 $A$ 得
\begin{equation}\label{eq:em-total-probability}
 \mathbb P(A)=\sum_{i=1}^m\mathbb P(B_i)\mathbb P(A\mid B_i).
\end{equation}
若 $\mathbb P(A)>0$，再用条件概率定义得到\emph{贝叶斯公式}（Bayes' formula）
\[
 \mathbb P(B_j\mid A)=
 \frac{\mathbb P(B_j)\mathbb P(A\mid B_j)}
 {\sum_i\mathbb P(B_i)\mathbb P(A\mid B_i)}.
\]
这是从观察结果更新来源概率的规则，而非交换条件符号。分母必须包含全部可能来源，条件概率也不是来源的先验概率。

\section{有限随机变量、期望与方差}
\emph{随机变量}（random variable）是函数 $X:\Omega\to\mathbb R$，$X$ 表示随机对象，$x$ 表示某个取值。其分布列为 $p_X(x)=\mathbb P(X=x)$，只需列出实际可能值。定义
\[
 \mathbb E[X]=\sum_x x\,p_X(x),\qquad
 \operatorname{Var}(X)=\mathbb E[(X-\mathbb E[X])^2]
 =\mathbb E[X^2]-\mathbb E[X]^2.
\]
所有和有限，故无可积性问题。按基本结果逐项相加可证明 $\mathbb E[aX+bY]=a\mathbb E[X]+b\mathbb E[Y]$，这不要求独立；独立时才可进一步保证乘积期望分解以及方差相加。

取值 $0,1$ 且 $\mathbb P(X=1)=p$ 的变量称为 Bernoulli 变量，它的期望为 $p$、方差为 $p(1-p)$。$n$ 次联合独立、成功概率同为 $p$ 的试验成功总数 $S$ 服从\emph{二项分布}（binomial distribution）。当 $0<p<1$ 时，
\[
 \mathbb P(S=k)=\binom nk p^k(1-p)^{n-k},\quad 0\le k\le n,
 \qquad \mathbb E[S]=np,\quad\operatorname{Var}(S)=np(1-p).
\]
概率式先指定哪 $k$ 次成功再求和；期望由线性性，方差由独立性推出。$p=0,1$ 时分别为恒等于 $0,n$ 的退化情形，无需给 $0^0$ 赋值。期望相同的变量可能具有不同方差，因此期望不完整描述结果波动。

\section{例题、习题与学习检查}
\begin{example}
等可能选择两个盒子之一，第一盒含两红一白，第二盒含一红三白。从选中盒中等可能取一球，观察为红球，求来自第一盒的概率。
\end{example}
\begin{solve}
两来源先验概率均为 $1/2$，红球条件概率分别为 $2/3,1/4$，故红球总概率为 $1/3+1/8=11/24$。第一盒且红球的概率为 $1/3$，后验概率为 $(1/3)/(11/24)=8/11$。
\end{solve}
\begin{enumerate}
\item 三次独立公平掷币恰有两次正面的概率是多少？答：$3/8$。
\item 公平六面骰点数的期望是多少？答：$7/2$；期望未必是可能取值。
\item 若 $Y=X$，其中 $X$ 是公平 Bernoulli 变量，比较 $\operatorname{Var}(X+Y)$ 与 $\operatorname{Var}(X)+\operatorname{Var}(Y)$。答：分别为 $1$ 与 $1/2$，缺少独立性时不能直接相加。
\end{enumerate}

